\chapter{Anlage starten und einrichten}
\label{ch:anlage}

Quelle der Schritte: \pfad{docs/inbetriebnahme.md} im Repository der Gruppe, ergänzt um die
Werkzeugdaten aus \pfad{code/m2_robot_control/ANLEITUNG_NEURA.md}.

\section{Vor dem Einschalten}

\begin{itemize}
  \item Zu zweit arbeiten, eine Person hat die Hand am Not-Halt.
  \item Arbeitsraum frei; Kamerastange, Prüfvorrichtung und Tisch an ihrem Platz.
  \item Handgelenkkamera fest am Greifer, Objektiv fixiert, Kabel so geführt, dass sie bei
    Drehung des Handgelenks nicht spannen.
  \item Kameralicht an der Handgelenkkamera einschalten (Standard seit 17.11.).
\end{itemize}

\section{Anlage einschalten}
\label{sec:einschalten}

Die Reihenfolge ist wichtig: Kameras und das Ethernet-Kabel des Roboters werden erst am Ende
an den Rechner gesteckt.

\begin{enumerate}
  \item \textbf{Hardware anschließen:} Roboter, Greifer usw. verbinden. Kameras und
    Roboter-Ethernet noch \emph{nicht} an den Rechner stecken.
  \item \textbf{Robotersteuerung einschalten} (Control Box).
  \item \textbf{Touchpad (Teach-Pendant) einschalten und anmelden.} Bleibt der Bildschirm
    schwarz: Kabel des Touchpads abziehen und wieder einstecken. Die Zugangsdaten gibt es beim
    Laborbetreuer; sie gehören nicht ins Repository.
  \item \textbf{Roboter einschalten:} Am Touchpad auf \textbf{Real} umschalten, dann den
    \textbf{weißen Knopf} an der Steuerung drücken $\rightarrow$ \textbf{Power On}. Falls nötig
    vorher den Not-Halt entriegeln und Fehler quittieren.
  \item \textbf{Greifer-Ausgang setzen:} \emph{Settings $\rightarrow$ IO Settings},
    \textbf{Output 1.3 auf 1}. Alle anderen Outputs müssen auf 0 stehen
    (\cref{sec:greifer}).
  \item \textbf{Kameras und Ethernet anstecken:} beide Daheng-Kameras direkt an getrennte
    USB-3-Anschlüsse, dann das Ethernet-Kabel des Roboters an den Rechner.
  \item \textbf{Oberfläche starten} (\cref{ch:oberflaeche}) und verbinden.
\end{enumerate}

\paragraph{Checkliste.}
\begin{itemize}
  \item[$\square$] Roboter, Greifer usw. angeschlossen
  \item[$\square$] Robotersteuerung eingeschaltet
  \item[$\square$] Touchpad an (bei schwarzem Bildschirm Kabel ab- und wieder anstecken), angemeldet
  \item[$\square$] auf Real geschaltet, weißer Knopf $\rightarrow$ Power On
  \item[$\square$] IO Settings: Output 1.3 = 1, alle anderen = 0
  \item[$\square$] Kameras angesteckt, Kameralicht an
  \item[$\square$] Roboter-Ethernet angesteckt
  \item[$\square$] Oberfläche gestartet, Anzeige rot nach dem Verbinden
\end{itemize}

\section{Greifer und Werkzeug einrichten}
\label{sec:greifer}

Montiert ist ein Zimmer LWR50L (elektrisch, IO-Link) auf dem Robotermodul LWR50F. In der
Steuerung ist dafür das Werkzeug \code{ZimmerLWR50} angelegt.

\begin{enumerate}
  \item Greifer mit Spannung versorgen: \emph{Settings $\rightarrow$ IO Settings},
    Output 1.3 = 1, alle anderen Outputs = 0 (Schritt 5 in \cref{sec:einschalten}).
  \item Am Pendant das Werkzeug \code{ZimmerLWR50} auswählen. Alle Punkte und alle Ziele von
    Modul~4 sind \acs{tcp}-Posen dieses Werkzeugs; mit einem anderen Werkzeug liegt jede
    Greifhöhe falsch.
  \item \acs{tcp}-Versatz prüfen: \SI{197.5}{\milli\metre} in $z$ (Flansch bis Backenspitzen,
    geschlossen).
  \item Last prüfen: Masse \SI{1.75}{\kilo\gram}, Schwerpunkt etwa \SI{81}{\milli\metre} über
    dem Flansch in $z$. Masse und Schwerpunkt sind wichtig für Schwerkraftkompensation und
    Kollisionserkennung.
  \item Greifer einmal öffnen und schließen (Pendant oder Reiter M2: \knopf{Greifer auf/zu}).
  \item In der Oberfläche von Modul~4 \emph{Systemcheck} $\rightarrow$
    \knopf{Steuerung pruefen (nur lesen)}: Greifermodus muss \enquote{hardware} sein, das
    Werkzeug gesetzt.
\end{enumerate}

\hinweis{Der Greifer kennt über NeuraPy nur \enquote{auf} und \enquote{zu}. Öffnungsweite und
Kraft sind feste Parameter des Werkzeugs in der Steuerung. Zwischen Bauteil
(\SI{39.7}{\milli\metre}) und Backen (\SI{44.8}{\milli\metre}) bleiben nur etwa
\SI{2.5}{\milli\metre} je Seite.}

\section{Punkte für Modul~4}
\label{sec:punkte}

Modul~4 liest die Punkte bei jedem Start direkt aus der Punktedatenbank der Steuerung.
\Cref{tab:punkte} zeigt, welche Punkte geteacht sein müssen.

\begin{table}[htbp]
  \centering
  \small
  \caption{Punkte in der Steuerung}
  \label{tab:punkte}
  \begin{tabularx}{\textwidth}{@{}l X@{}}
    \toprule
    Punkt & Bedeutung \\
    \midrule
    \code{CLEAR_FOV} & Startstellung; der Arm verdeckt die Draufsichtkamera nicht \\
    \code{PICK} & Greifpunkt am Bauteil; bei automatischen Lagen setzt die Aufnahme ihn selbst \\
    \code{PRE_PLACE} & Übergabepunkt über der Prüfvorrichtung (nur Ablauf
      \pfad{pick_to_station}) \\
    \code{ID0} bis \code{ID8} & Tischpunkte mittig über den ArUco-Markern im $3\times3$-Raster,
      Backen berühren den Tisch \\
    \code{IGNORE} & Mitte einer Sperrzone (Größe in der Ablaufdatei) \\
    \bottomrule
  \end{tabularx}
\end{table}

\code{ANFLUG}, \code{PRE_PICK} und der Endpunkt \code{UEBER_TISCH} werden aus \code{PICK} bzw.\
\code{ID4} berechnet und müssen nicht geteacht werden.

\achtung{Werden die Tischpunkte neu geteacht, ändert sich das Kollisionsmodell. Die Punkte vor
jeder Änderung sichern: \emph{Systemcheck} $\rightarrow$ \knopf{Punkte sichern} schreibt eine
Datei nach \pfad{sequences/}.}

\section{Erste Prüfung nach dem Einschalten}

In der Oberfläche von Modul~4, Reiter \emph{Systemcheck}, Gruppe \enquote{verbunden, ohne
Bewegung}:
\begin{enumerate}
  \item \knopf{Steuerung pruefen (nur lesen)}: \code{is_robot_in_simulation() = False},
    Werkzeug, Greifermodus.
  \item \knopf{Kinematik und Achsgrenzen}: Abweichung zur \acs{vm} unter
    \SI{1}{\milli\metre}.
  \item \knopf{Kameras suchen} und \knopf{Kameras: Rate und Versatz}: beide Daheng mit
    Seriennummer, je mindestens \SI{15}{fps}, Bildalter p95 unter \SI{30}{\milli\second}.
\end{enumerate}
Danach \knopf{Protokoll speichern}; das Ergebnis liegt in \pfad{Berichte/<Datum>_Systemcheck.txt}.
